\chapter{STUDY READINESS AND SOURCE REGISTER}
\label{app:readiness}
\newpage
\section{Implemented, Proposed, and Unresolved}
The source snapshot documents one practice trial in the first assigned voice, two 40-trial measured blocks, nominal-angle quotas, right-trigger selection, and a final external questionnaire handoff. The protocol proposes exact pre-cue geometry and balanced cross-relative heights; the October 20 audit identifies those as implementation gaps. These statuses must be rechecked against the final release rather than inferred from the presence of future-tense method text.

\begin{longtable}{@{}p{0.27\linewidth}p{0.64\linewidth}@{}}
\caption{Decisions required before confirmatory collection.}\\\toprule
Requirement & Completion evidence \\\midrule\endfirsthead
\toprule Requirement & Completion evidence \\\midrule\endhead
Starting geometry & Frozen origin/cross/target snapshot before cue, signed/absolute angle, time and frame provenance.\\
Vertical allocation & Seated-safe bounds and near-zero rule; realized 20-above/20-below schedule per block.\\
Power and model & Meaningful effect, theta function, random effects, simulation assumptions, participant N and enrollment ceiling.\\
Tracking acceptance & Independent spatial validation, unique-source timing, drift checks, and stop/recalibration rules.\\
Practice proficiency & Observable criterion and response if a single practice trial is insufficient.\\
Timeout and stopping & Prespecified comfort, safety, technical failure, and incomplete-trial treatment.\\
Acoustic checks & Per-voice intelligibility, comfortable levels, wording verification, and actual delivery timing.\\
Survey readiness & Burden/comprehension pilot, final retrospective references, validated scoring/export implementation.\\
Consent & Actual cloud processing, identifiable voice handling, storage, deletion, and approved participant wording.\\
Release freeze & Build, policy, phrase, provider and instrument versions; full-session wearer acceptance.\\\bottomrule
\end{longtable}

\section{Source Register}
The current Overleaf project is \url{https://www.overleaf.com/project/6a84bc097073254fbca6985b}. This first pass uses revision \texttt{f9f49e450c7b64b34c319a557114f1fa4393f003}, retrieved without overwriting the working Overleaf project. The system description preserves the current Section 3 as its foundation, with thesis-level expansion and language edits. The method and planned-analysis chapters similarly expand the current standalone sections. The 62-entry bibliography is retained as the working library; only cited entries appear in the reference list.

The full protocol handoff dated October 20 supersedes the older handoff on practice count and survey timing. The later user clarification assigns practice to the first counterbalanced voice. The repository alignment audit supplies implementation status and recorded test evidence. Historical four-condition, 14-round, dwell-selection, and immediate post-block questionnaire descriptions are not the design for this thesis.

The supplied UCF LuaLaTeX template provides the class and accessibility setup. The downloaded reference is Hoang Bao Truong's \emph{Using Inferences from Natural Language Narratives to Improve Video Question Answering}, UCF, 2026, available at \url{https://stars.library.ucf.edu/gradstudies_etd_2026/199/}. Its 140-page repository PDF includes a repository cover sheet, front matter, detailed methodology, results, conclusion, appendices, and references. It is a presentation reference, not a source of findings or text for this project.

\section{Figure Provenance}
The scene illustration, procedure diagram, and guidance-memory architecture diagram were already present in the current Overleaf project and are reused here. Their editable SVG sources are included where available. Ten additional vector diagrams explain the comparison, state families, memory, geometry, event boundaries, data flow, and analysis scope. An actual Trial Review capture shows its explicitly synthetic bundled demonstration. Figure sources and provenance are listed in \path{figures/README.md}. These figures are explanatory schematics or interface illustrations, not measured participant results. Drawn object counts and highlight colors must not be mistaken for the frozen participant release or for an identifying target cue presented before search. The current audit describes 168 spawned objects, 21 per wall; the eight trials per nominal theta category are a different quantity.

\section{Submission Work Remaining}
This draft requires author and advisor review, a verified final study record, completed empirical analyses, and updated conclusions. Title-page administrative details and the final AI disclosure must be confirmed. Citation metadata and source-to-claim fit should receive a final audit, including any newly added quantitative literature claims. Accessibility tagging, reading order, table headers, and figure descriptions require final validation with the university's submission process. Successful compilation alone is not certification of ETD compliance or accessibility.
